\par\noindent\SourceLeaf{51}%
\hspace*{\parindent}6. Nun können wir ohne weiteren Aufschub zum gesteckten Ziel übergehen. Man entwickle $(1+n\cos\varphi)^{\nu}$ nach dem binomischen Lehrsatz in eine Reihe nach Potenzen von $\cos\varphi$. Sodann löse man jedes Glied dieser Reihe in eine neue Reihe nach Kosinus von Vielfachen des gegebenen Winkels auf. Aus diesen Reihen wähle man die Glieder aus, die mit dem Kosinus desselben Winkels multipliziert sind, und fasse sie zu einem einzigen Glied zusammen. Kennt man das Gesetz der auf diese Weise entstehenden Glieder, so ist die ganze Aufgabe nahezu gelöst. In der Grundreihe

\[
\begin{aligned}
(1+n\cos\varphi)^{\nu} & = 1 + {}^{\nu}\mathfrak{B}^{(1} n\cos\varphi + {}^{\nu}\mathfrak{B}^{(2)} n^{2}\cos\varphi^{2} \\&\quad + {}^{\nu}\mathfrak{B}^{(3)} n^{3}\cos\varphi^{3} + .\,. + {}^{\nu}\mathfrak{B}^{(q)} n^{q}\cos\varphi^{q} + .\,.
\end{aligned}
\]

\noindent kommen der Reihe nach alle Potenzen von $\cos\varphi$ vor, die geraden wie die ungeraden. Aus den geraden entstehen bei der Entwicklung nur Kosinus gerader Vielfacher, aus den ungeraden nur solche ungerader Vielfacher. Geht es also um den Kosinus eines ungeraden Vielfachen, so muss man auf die ungeraden Potenzen zurückgreifen, geht es um den Kosinus eines geraden Vielfachen, auf die geraden.

A) Man erhält das Anfangsglied der gesuchten Reihe, indem man aus allen vorgegebenen Gliedern, deren Entwicklungen mit $\cos 0\varphi = 1$ beginnen, die Anfangsglieder in ihrer ursprünglichen Reihenfolge vereinigt. Nur gerade Potenzen von $\cos\varphi$ liefern bei der Entwicklung ein solches Anfangsglied. Jede von ihnen liefert einen hierher gehörigen Teil. Der Teil $h^{\text{ta}}$ geht also aus dem Glied $2h^{\text{to}}$ der Grundreihe hervor, das ${}^{\nu}\mathfrak{B}^{(2h)} n^{2h}\cos\varphi^{2h}$ ist. Das Anfangsglied von $\cos\varphi^{2h}$, das allein hierher gehört, ist

\[
= \frac{1}{2^{2h-1}} \frac{{}^{2h}\mathfrak{B}^{(h)}}{2} = \frac{1}{2^{2h}} {}^{2h}\mathfrak{B}^{(h)};
\]

\noindent und unter Hinzunahme des Koeffizienten von $\cos\varphi^{2h}$ erhält man damit den Teil $h^{\text{ta}}$ der mit $\cos 0\varphi = 1$ multiplizierten Reihe

\[
\frac{1}{2^{2h}} {}^{2h}\mathfrak{B}^{(h)}  .  {}^{\nu}\mathfrak{B}^{(2h)} n^{2h}.
\]

\noindent So ist das allgemeine Glied $h^{\text{tus}}$ der Reihe gefunden, die das Anfangsglied jener Reihe bildet, die man durch Entwicklung von $(1+n\cos\varphi)^{\nu}$ nach den Kosinus von Vielfachen des Winkels erhält. Dass ihr Anfangsglied $=1$ ist, liegt auf der Hand. Setzen wir nun

\[
(1+n\cos\varphi)^{\nu} = A + A^{(1)}\cos\varphi + A^{(2)}\cos 2\varphi + .\,.\,.,
\]

\noindent so ergibt sich

\[
A = 1 + \frac{1}{2^{2}} {}^{2}\mathfrak{B}^{(1)}  .  {}^{\nu}\mathfrak{B}^{(2)} n^{2} + \frac{1}{2^{4}} {}^{4}\mathfrak{B}^{(2)}  .  {}^{\nu}\mathfrak{B}^{(4)} n^{4} + .\,.\,. + \frac{1}{2^{2h}} {}^{2h}\mathfrak{B}^{(h)}  .  {}^{\nu}\mathfrak{B}^{(2h)} n^{2h} + .\,.\,.
\]

\noindent oder

\[
A = 1 + \frac{1}{2^{2}}  .  \frac{2}{1} \frac{\nu.\nu-1}{1.2} n^{2} + \frac{1}{2^{4}}  .  \frac{4.3}{1.2}  .  \frac{\nu.\nu-1.\nu-2.\nu-3}{1.2.3.4} n^{4} + \text{etc.}
\]
